\newpage
\section{Einleitung (Ronja Rüther)}
Digitale Sprachlernanwendungen gewinnen in den letzten Jahren stetig an Bedeutung
zum selbstständigen Erlernen einer Fremdsprache. Da die meisten Personen heutzutage ein
Smartphone verfügbar haben, lassen sich Vokabeln regelmäßig wiederholen und der Wortschatz
festigen, da Wartezeiten so überall zum Lernen genutzt werden können.
Die bekannte Anwendung \textit{Duolingo} zeigt bereits, wie das digitale Erlernen einer neuen Sprache 
funktionieren kann. Um das Lernen einfacher und motivierender zu gestalten, werden 
Gamification-Elemente, wie z.B. das Zusammensetzen eines Satzes aus Wort-Bausteinen, verwendet. \parencite{Duolingo}
Dennoch stößt diese Plattform an eine Grenze, denn die Lerninhalte sind statisch vorgegeben
und der Nutzer kann keine eigenen Vokabeln anlegen, die er speziell für die Schule, den Beruf oder ein
neues Hobby benötigt.

Ziel dieses Projekts ist daher die Entwicklung einer Webanwendung namens \enquote{Lingotrain},
die das Prinzip der interaktiven Sprachübungen mit der eigenen Verwaltung der Lerninhalte kombiniert.
Statt vorgefertigter Lektionen ermöglicht diese Anwendung die Erstellung individueller Lernsets mit Vokabeln.
Der Fokus liegt dabei auf den Zielsprachen Englisch, Französisch und Spanisch mit Deutsch als Ausgangssprache.
Durch dynamisch generierte Übungen und Gamification-Aspekte können die Nutzer ihre eigenen 
Vokabeln im eigenen Tempo lernen.
Damit die Vokabeln, insbesondere einzelne Wörter, auch kontextbezogen und im Hinblick auf eine 
Unterhaltung gelernt werden, verfügt die Webanwendung auch über einen Chatbot, der passende Sätze 
zu den einzelnen Wörtern generiert und dem Nutzer passendes Feedback geben kann. Darüber hinaus ist auch
ein Konversationsmodus geplant, in dem der Nutzer mit dem Chatbot eine Unterhaltung zu einer bestimmten
Vokabel aus seinen Lernsets führen kann.

Die technische Umsetzung von \enquote{Lingotrain} stellt dabei hohe Anforderungen an die Architektur
des verwendeten Frontend-Frameworks. Da neben statischen Vokabellisten auch reaktive Übungsmodule sowie
asynchrone Chat-Schnittstellen nötig sind, wird eine modulare und komponentenbasierte Struktur benötigt.
Im Rahmen dieser Ausarbeitung wird daher dokumentiert, wie die Anforderungen mithilfe eines modernen und 
komponentenorientierten Frameworks umgesetzt wurden, damit der Anwender die Webanwendung performant nutzen 
kann.

\section{Verwendete Bibliotheken und Werkzeuge}

\section{Aufbau der Webanwendung}

\section{Architektur}

\section{Aufgetretene Probleme}

\section{Ergebnisse und Fazit}

Eine Aufzählung:

\begin{itemize}
\item Lorem ipsum
\item dolor sit amet
\item consetetur sadipscing elitr
\item sed diam nonumy
\end{itemize}
Eine nummerierte Aufzählung:

\begin{enumerate}
\item Erster Punkt
\item Zweiter Punkt
\item Dritter Punkt
\end{enumerate}
Es folgt Tabelle \ref{beispieltabelle}.

\begin{table}[htbp]
\centering
\begin{tabular}{lrc}
\toprule
Linksbündig & Rechtsbündig & Zentriert \\
\midrule
Lorem &  13 & amet \\ \addlinespace
ipsum & 104 & consetetur \\ \addlinespace
dolor &   7 & sadipscing \\ \addlinespace
sit   &  -5 & elitr \\
\bottomrule
\end{tabular}
\caption{Beschreibung der Tabelle.}
\label{beispieltabelle}
\end{table}

Programmcode im Fließtext: \lstinline{printf("Hello, world!\n");}.
Lorem ipsum dolor sit amet, consetetur sadipscing elitr, sed diam nonumy eirmod tempor invidunt ut labore et dolore magna aliquyam erat, sed diam voluptua. At vero eos et accusam et justo duo dolores et ea rebum. Stet clita kasd gubergren, no sea takimata sanctus est Lorem ipsum dolor sit amet.
Es folgt Programmlisting \ref{beispiellisting}.

\begin{lstlisting}[caption={Beschreibung des Listings.}, label=beispiellisting, frame=tblr, numbers=left]
#include <stdio.h>

int main(void)
{
    printf("Hello, world!\n");
}
\end{lstlisting}
Ein Listing ohne Titel, welches nicht im Listingverzeichnis aufgefühert wird:

\begin{lstlisting}
#include <stdio.h>

int main(void)
{
    printf("Hello, world\n");
}
\end{lstlisting}

Text kann \emph{kursiv} oder \textbf{fett} gesetzt werden.
Lorem ipsum dolor sit amet, consetetur sadipscing elitr, sed diam nonumy eirmod tempor invidunt ut labore et dolore magna aliquyam erat, sed diam voluptua. At vero eos et accusam et justo duo dolores et ea rebum. Stet clita kasd gubergren, no sea takimata sanctus est Lorem ipsum dolor sit amet. Lorem ipsum dolor sit amet, consetetur sadipscing elitr, sed diam nonumy eirmod tempor invidunt ut labore et dolore magna aliquyam erat, sed diam voluptua. At vero eos et accusam et justo duo dolores et ea rebum. Stet clita kasd gubergren, no sea takimata sanctus est Lorem ipsum dolor sit amet.
